\begin{proposition}[toy-RH 临界线模态的 \texorpdfstring{$2$}{2}-torsion 谱来源]\label{prop:pom-toy-rh-2torsion-source}
令 \(M\) 为真实输入 \(40\) 状态核的邻接矩阵。
在命题 \ref{prop:pom-zeta-factorization-40} 的分解下，
\[
\zeta_M(z)
=
\frac{1}{(1-z)^2(1+z)}\cdot \zeta_{A\otimes A}(z)\cdot \zeta_{-A}(z),
\qquad
\zeta_{-A}(z)=\frac{1}{\det(I-z(-A))}.
\]
则 Dirichlet 坐标 \(z=\lambda^{-s}\)（其中 \(\lambda=\rho(M)=\varphi^2\)）下，落入开带 \(0<\Re(s)<1\) 的全部非平凡极点列必完全来自 \(\zeta_{-A}\) 因子，并且其相位等价于阶 \(2\) 的符号扭曲（\(\pm 1\) 取值的角色）所贡献的通道模态。
因此该临界线晶格模态可被解释为一条 \(\ZZ/2\) 型 torsion 的扭曲谱模态；在 \v{C}ech 语言中，对应于系数推前后取值为 \((-1)\) 的单位根相位缺陷。
\end{proposition}

\begin{proof}
命题 \ref{prop:pom-zeta-factorization-40} 与注 \ref{rem:pom-minusA-factor-critical-line} 已将 \(\zeta_M\) 的结构性分解与“临界线晶格”的来源分离出来：\(-A\) 的 Perron 根为 \(-\varphi=-\sqrt{\lambda}\)，其在 \(z=\lambda^{-s}\) 下对应 \(|\lambda z|=|\lambda^{1-s}|=\sqrt{\lambda}\)，因而给出 \(\Re(s)=\tfrac12\) 的竖直晶格（定理 \ref{thm:pom-toy-rh}）。
另一方面，\(\zeta_{A\otimes A}\) 的非平凡谱点模长为 \(1\) 或 \(\lambda^{-1}\)，对应的 \(s\)-像落在 \(\Re(s)\in\{1,0\}\) 的边界线上；显式短周期因子 \((1-z)^2(1+z)^{-1}\) 亦不产生开带内的新极点。
因此，开带内的全部非平凡极点列只能来自 \(\zeta_{-A}\) 的符号扭曲模态。
由于该扭曲仅改变符号，其角色必为阶 \(2\)（单位根为 \(-1\)），从而得到所述 \(2\)-torsion 谱来源解释。证毕。
\end{proof}

\endinput
